\subsection{Financial Literacy and Personal Financial Decision-Making}
Financial literacy is widely treated as a foundation for informed household financial decisions. Lusardi and Mitchell synthesize evidence linking financial knowledge with decisions involving saving, borrowing, investment, and retirement planning, while also identifying implications for education and public policy \cite{lusardiMitchell2014,lusardiMitchell2023}. This body of work motivates the inclusion of structured learning material in a financial decision-support system. It does not, by itself, establish that a digital learning module produces measurable improvements in users' knowledge; such an effect requires evaluation with learners.

Emergency savings and financial resilience form a complementary theme. Research on household financial capability has examined how knowledge, confidence, and access to savings relate to preparedness for unexpected expenses. These findings motivate the use of cash-flow and emergency-coverage indicators in personal finance tools, while cautioning against treating any single coverage threshold as universally appropriate.

\subsection{Robo-Advisory and Automated Investment Support}
Robo-advisory research examines automated investor profiling, portfolio construction, investor behavior, service performance, and adoption. Cardillo and Chiappini's systematic review identifies multiple streams in the robo-advisor literature, including behavioral research, performance, and algorithmic modelization \cite{cardilloChiappini2024}. Such work provides a useful foundation for automated investment decision support, but robo-advisory systems differ in their investment universe, investor data, regulatory context, and underlying recommendation logic.

For an Indian retail-investor context, automated support should also account for users' understanding of investment concepts and the interpretability of recommendations. Research on robo-advisor adoption and user experience has documented transparency and comprehension as important concerns. Zhu et al.'s qualitative study of users interacting with an AI-enabled financial advisory system reported that participants found the system insufficiently transparent and information difficult to understand \cite{zhuEtAl2023}. This supports a design emphasis on exposing recommendation constraints and reasons rather than presenting only an output ranking.

\subsection{Explainability in Financial Decision Systems}
Explainable artificial intelligence has received increasing attention in finance, where decisions may need to be traceable and understandable. Weber, Carl, and Hinz reviewed 2,022 candidate articles and classified 60 relevant studies across finance, information systems, and computer science, describing explainability methods and their applications in areas such as risk management and portfolio optimization \cite{weberCarlHinz2024}. Digmayer's study of robo-advisor adoption barriers likewise highlights user needs related to understandability and trustworthiness across configuration, matching, and maintenance stages \cite{digmayer2024}.

The literature does not justify assuming that explanations automatically increase trust, adoption, or recommendation adherence. Accordingly, WealthGenie uses deterministic rule-based explanations to disclose profile factors, hard eligibility constraints, scoring dimensions, and scenario assumptions. In this paper, explainability is considered a system property—whether the decision path and its rationale are exposed—not a demonstrated behavioral outcome.

\subsection{Positioning and Research Gap}
The reviewed work addresses financial literacy, automated investment support, and explainability, but these strands do not necessarily appear together in one educational prototype. WealthGenie integrates structured financial-literacy lessons, financial readiness assessment, a rule-based mutual-fund recommendation workflow, paper trading and historical replay, and a scenario-based financial Digital Twin in an Indian personal-finance setting. The contribution is therefore the design and implementation of an integrated, transparent decision-support prototype, together with deterministic and reproducibility-oriented technical evaluation.

The current evaluation is limited to software-level checks and generated profiles. It does not include a real-user learning study, expert suitability assessment, or evidence of improved investment outcomes. These questions remain for future empirical work.
